\section{Introduction}

The emergence of decentralized finance represents one of the most significant structural transformations in financial intermediation since the advent of electronic trading. Between January 2020 and December 2025, the total value locked (TVL) in DeFi protocols expanded from under \$1 billion to over \$100 billion at its peak, with lending protocols consistently accounting for the largest share of this activity~\cite{schar2021decentralized2, jiang2023defi}. Unlike traditional financial intermediaries, which rely on institutional credit assessment, relationship-based lending, and regulatory compliance frameworks to manage risk, DeFi lending protocols operate through immutable smart contracts that algorithmically enforce collateralization requirements, interest rate adjustments, and automatic liquidation of under-collateralized positions~\cite{gudgeon2020defi}. The complete transparency of these protocols means that every deposit, withdrawal, borrowing event, repayment, collateral adjustment, and liquidation is immutably recorded on a public blockchain, creating an unprecedented empirical window into the behavior of financial market participants operating under varying degrees of risk exposure.

This transparency has given rise to a rapidly growing literature that examines DeFi lending from multiple analytical angles. One strand of research investigates the systemic risks inherent in these protocols, including liquidity crunches, cascading liquidations, and the potential for contagion across composable DeFi primitives~\cite{qin2021empirical, perez2020liquidations}. A second strand develops risk modeling frameworks tailored to the unique features of DeFi lending, including multi-chain environments, cross-protocol dependencies, and institutional-grade risk taxonomies~\cite{oberholzer2026toward, sevim2026interoperability}. A third strand applies machine learning methods to on-chain data for credit scoring and default prediction~\cite{ghosh2024onchain}. A fourth strand examines the engineering properties of specific protocol mechanisms, including liquidation thresholds, interest rate curves, and governance design~\cite{iftikhar2025automated, tovanich2023contagion}. 

Despite the richness of this literature, a notable gap persists at its core. The existing research on DeFi lending risk has focused overwhelmingly on what might be termed \emph{state variables}—static snapshots of position health captured by health factors, collateralization ratios, and borrowing volumes at discrete time points. What has received considerably less empirical attention is the \emph{behavioral process} through which borrowers navigate rising position risk: the active decisions they make to adjust collateral, repay debt, increase their exposure, or maintain their positions as their risk profile deteriorates. This distinction matters because two borrowers with identical health factors at a given moment may follow fundamentally different trajectories depending on their behavioral responses to mounting risk pressure. One borrower may proactively add collateral and reduce leverage early in the risk cycle, while another may maintain or even increase exposure until forced into liquidation by third-party arbitrageurs. Conventional risk indicators would place these two borrowers at the same level of risk, yet their behavioral processes suggest meaningfully different approaches to risk management that may contain information about future outcomes.

This research report addresses this gap by proposing a study situated at the intersection of DeFi lending infrastructure, behavioral finance theory, and on-chain credit risk assessment. The central research question we propose is: In public on-chain lending markets, do borrowers' active adjustments under rising position risk provide risk information beyond that contained in conventional on-chain indicators? This question decomposes into two analytically distinct but empirically connected layers. The first, behavioral layer asks whether borrowers exhibit systematic patterns of active adjustment as their positions approach risky thresholds—patterns that are distinguishable from random or purely mechanical responses to price movements. The second, credit-oriented layer asks whether these behavioral patterns, if they exist, carry incremental predictive power for future liquidation events or risk deterioration beyond what is already captured by health factors, collateral ratios, and account activity measures.

%
% Figure 1: Conceptual Framework
%
\begin{figure}[t]
\centering
\resizebox{\textwidth}{!}{%
\begin{tikzpicture}[
  node distance=1.2cm and 0.6cm,
  box/.style={rectangle, rounded corners=2pt, draw=#1!60, fill=#1!12, text=#1!80!black, minimum width=2.2cm, minimum height=0.65cm, align=center, font=\scriptsize\sffamily, line width=0.4pt},
  widebox/.style={rectangle, rounded corners=2pt, draw=#1!60, fill=#1!12, text=#1!80!black, minimum width=9cm, minimum height=0.65cm, align=center, font=\small\sffamily\bfseries, line width=0.4pt},
  smallbox/.style={rectangle, rounded corners=2pt, draw=#1!60, fill=#1!10, text=#1!80!black, minimum width=1.8cm, minimum height=0.5cm, align=center, font=\tiny\sffamily, line width=0.3pt},
  arrow/.style={->, >=stealth, line width=0.6pt, draw=black!55},
  darrow/.style={<->, >=stealth, line width=0.7pt, draw=black!55},
]
  % Layer 1: Protocols
  \node[widebox=blue, label={[font=\footnotesize\sffamily\bfseries]north:Data Sources}] (proto) {Aave (V2/V3) \hspace{0.8cm} Compound (V2/V3) \hspace{0.8cm} MakerDAO};

  % Layer 2: Behavioral variables  
  \node[widebox=teal, below=0.7cm of proto, label={[font=\footnotesize\sffamily\bfseries]north:Behavioral Process Extraction}] (behave) {\normalsize Borrower Active Adjustment \quad (collateral changes, debt management, response timing)};

  % Layer 3: Two tracks
  \node[box=violet, below left=1.1cm and -1.2cm of behave, minimum width=3.8cm, font=\scriptsize\sffamily] (rq1) {\textbf{RQ1: Behavioral Layer}\\[2pt]Do borrowers systematically adjust\\as HF $\rightarrow$ 1.0?};
  \node[box=violet, below right=1.1cm and -1.2cm of behave, minimum width=3.8cm, font=\scriptsize\sffamily] (rq2) {\textbf{RQ2: Credit Layer}\\[2pt]Does adjustment behavior predict\\future liquidations?};

  \node[box=red!80!black, below=0.6cm of rq1, minimum width=3.8cm, font=\scriptsize\sffamily] (out1) {Prospect Theory Validation\\(loss aversion, reference-point)};
  \node[box=red!80!black, below=0.6cm of rq2, minimum width=3.8cm, font=\scriptsize\sffamily] (out2) {Behavioral Early Warning\\(credit risk signals)};

  % Arrows
  \draw[arrow, thick] (proto.south) -- (behave.north);
  \draw[arrow, thick] (behave.south) |- (rq1.north);
  \draw[arrow, thick] (behave.south) |- (rq2.north);
  \draw[arrow] (rq1.south) -- (out1.north);
  \draw[arrow] (rq2.south) -- (out2.north);
  \draw[darrow, thick] (rq1.east) -- node[above, font=\tiny\sffamily]{H2a: Incremental predictive power?} (rq2.west);

  % Left labels
  \node[font=\footnotesize\sffamily, text=black!60, rotate=90, above left=0.0cm and -1.8cm of behave, anchor=south] {$\longleftarrow$ Behavioral Finance $\longrightarrow$};
\end{tikzpicture}%
}
\caption{Conceptual framework for studying borrower behavior and credit signals in DeFi lending. The framework proceeds from raw on-chain data (top) through behavioral process variable extraction (center) to the research questions (bottom). RQ1 characterizes how borrowers adjust positions as risk accumulates; RQ2 tests whether these behavioral patterns provide incremental predictive power for future liquidation events beyond conventional risk indicators.}
\label{fig:framework}
\end{figure}

\vspace{4pt}

The report is organized as follows. Section~2 surveys the relevant literature across three domains: DeFi lending infrastructure and risk, behavioral finance and prospect theory, and on-chain credit risk assessment. Section~3 identifies the specific research topic and articulates the research questions and hypotheses. Section~4 describes the proposed research methodology, including data sources, variable construction, and analytical framework. Section~5 discusses possible outcomes, contributions, limitations, and boundary conditions.
